\section{The tangent index and restrictions on smooth fibres}

\subsection{The Picard family and the tangent index (C.1)}

\begin{proposition}[C.1.1; Proved]\label{st:C.1.1}
Let $Y$ be a connected compact complex threefold with
$b_1(Y)=b_2(Y)=0$. Every holomorphic line bundle $L$ on $Y$ satisfies
\[
 \chi(Y,T_Y\otimes L)=\frac12\int_Yc_3(T_Y)
 =1-\frac{b_3(Y)}2.
\]
The orientation in the integral is the complex orientation. The group
$\Pic(Y)$ has finitely many connected components, each a torsor for the
complex vector group $H^1(Y,\OO_Y)$, and
$\pi_0\Pic(Y)\simeq H^2(Y;\Z)$ through the first Chern class.
\end{proposition}
\begin{proof}
The universal coefficient theorem gives
$H^1(Y;\Z)=\Hom(H_1(Y;\Z),\Z)=0$ from $b_1(Y)=0$.
The group $H^2(Y;\Z)$ is finite because $Y$ has a finite CW model and
$b_2(Y)=0$. In the exponential sequence, normalized by
$\exp(2\pi i\,\cdot)$, the relevant maps are
\[
0=H^1(Y;\Z)\longrightarrow H^1(Y,\OO_Y)
 \xrightarrow{\exp} \Pic(Y)
 \xrightarrow{c_1} H^2(Y;\Z).
\]
They are group homomorphisms; the middle map is holomorphic on the
Picard group. Consequently its image $\Pic^0(Y)$ is isomorphic to the
complex vector space $H^1(Y,\OO_Y)$. The next map in the exponential
sequence is $H^2(Y;\Z)\to H^2(Y,\OO_Y)$; it is zero because its
source is finite and its target is a complex vector space. Thus
$c_1$ is surjective, and the component group is exactly $H^2(Y;\Z)$.

Poincar\'e duality gives $H^4(Y;\Q)=0$. Thus
$c_1(T_Y)$, $c_2(T_Y)$ and $c_1(L)$ vanish in rational cohomology.
The degree-six component of
$\operatorname{ch}(T_Y)\exp(c_1(L))\operatorname{td}(T_Y)$ is then
$c_3(T_Y)/2$. The Hirzebruch--Riemann--Roch formula for the Dolbeault
index gives the first equality. Chern--Gauss--Bonnet identifies
$\int_Yc_3(T_Y)$ with $e(Y)$. Since $Y$ is a connected closed oriented
six-manifold, its Betti numbers satisfy
$b_0=b_6=1$, $b_1=b_5=0$, and $b_2=b_4=0$. Hence
$e(Y)=2-b_3(Y)$. The pairing on $H^3(Y;\Q)$ is nondegenerate and
alternating, so $b_3(Y)$ is even, consistently with the integral index.
The exact symbolic expansion is recorded in
\texttt{code/D\_exact\_checks.py} and its output.
\end{proof}

\begin{proposition}[C.1.2; Proved]\label{st:C.1.2}
Let $Y$ be as in Statement C.1.1 and let $D$ be any nonzero effective
Cartier divisor, with arbitrary positive multiplicities. For every
holomorphic vector bundle $E$ and every component $P$ of $\Pic(Y)$ there
is a dense open subset $P_E\subset P$ on which
\[
 H^0(Y,E\otimes L)=H^3(Y,E\otimes L)=0.
\]
Each excluded locus is a proper closed complex analytic subset of $P$.
\end{proposition}
\begin{proof}
Choose $m>0$ annihilating $c_1(\OO_Y(D))$ in the finite group
$H^2(Y;\Z)$. Thus $\OO_Y(mD)\in\Pic^0(Y)$. Fix a representative $L_P$
of the component $P$. The inclusions of locally free sheaves
\[
 E\otimes L_P(-(k+1)mD)\hookrightarrow E\otimes L_P(-kmD)
 \hookrightarrow E\otimes L_P
\]
give a descending sequence of finite-dimensional spaces of global
sections. Its eventual stable value is zero. Indeed every section in
that value is locally divisible by every power of an equation of
$mD$. At a smooth point of a component of $D_{\rm red}$, the equation
has positive vanishing order. The intersection of all its powers in
the local ring of the manifold is zero. The section therefore vanishes
on a nonempty open set and hence on connected $Y$. We have proved
$H^0(Y,E\otimes L_P(-kmD))=0$ for some integer $k$.

We specify the family used in applying semicontinuity. Put
$V=H^1(Y,\OO_Y)$ and choose additive holomorphic \v Cech
one-cocycles $a^1_{ij},\ldots,a^q_{ij}$ representing a complex basis
on a common Stein open cover. On the product of this cover with $V$
use the transition functions
\[
 g_{ij}(v)=\exp\left(2\pi i\sum_{k=1}^q v_k a^k_{ij}\right).
\]
The additive cocycle identities imply $g_{ij}g_{jk}=g_{ik}$,
so these functions define a holomorphic line bundle $\mathscr P$.
Its restriction to $Y\times\{v\}$ is the line bundle
corresponding to $v$ under the preceding exponential convention.
Tensoring with $\operatorname{pr}_Y^*L_P$ gives a holomorphic family
parametrizing the entire component $P$. The projection
$Y\times V\to V$ is proper and the sheaf of sections of this bundle
is flat over $V$.

The proper cohomology-and-base-change theorem represents its cohomology
locally on $V$ by a finite complex of finite free analytic modules.
The conditions $h^j\geq a$ are determinantal analytic conditions on
the differentials of that complex. Thus the degree-zero jump locus is
closed analytic, and the line bundle $L_P(-kmD)$ constructed above
shows that it is proper. Its complement is dense because $V$ is
irreducible. Applying this argument to $E^*\otimes K_Y$ in the
component $P^{-1}$, and applying Serre duality
\[
 H^3(Y,E\otimes L)^*
 \simeq H^0(Y,E^*\otimes K_Y\otimes L^{-1}),
\]
gives the degree-three assertion. The intersection of the two dense
open subsets is dense open. The argument uses only positive
multiplicity at one component of $D$; no reduction of its scheme
structure was made.
\end{proof}



\begin{corollary}[C.1.3; Proved]\label{st:C.1.3}
For every $Y$ and $D$ in Statement C.1.2 and every $L\in\Pic(Y)$,
\[
 h^2(Y,T_Y\otimes L)\geq
 \max\left\{0,1-\frac{b_3(Y)}2\right\}.
\]
In particular, $b_3(Y)=0$ gives $H^2(Y,T_Y\otimes L)\ne0$ for
every $L$. A threefold in this class possessing a twist with vanishing
second tangent cohomology has $b_3(Y)\geq2$.
\end{corollary}
\begin{proof}
On the dense open subset of each Picard component supplied by
Statement C.1.2, Statement C.1.1 becomes
\[
 h^2(Y,T_Y\otimes L)-h^1(Y,T_Y\otimes L)
 =1-b_3(Y)/2.
\]
The asserted lower bound holds there. The jump locus for the lower
bound on $h^2$ is closed analytic in the explicit family above and
contains a dense subset of the component. It therefore contains the
entire component. This proves the bound for all line bundles,
including torsion bundles and the trivial bundle. The final assertion
is its direct numerical consequence.
\end{proof}

Statements C.1.1--C.1.3 prove Lead 4 with rational Betti-number
vanishing. No integral trivialization of either Chern class is needed
for the index. The existence of the divisor enters in the endpoint
vanishing, and a positive multiple of its actual Cartier divisor
enters the Picard family. These statements supply an index
obstruction; they make no assertion that the failed degree-two
vanishing holds.

\subsection{What follows for the smooth fibre (C.3)}

\begin{proposition}[C.3.1; Proved]\label{st:C.3.1}
Let $f:Y\to C$ be a holomorphic algebraic reduction of a connected
compact complex threefold with $b_1(Y)=b_2(Y)=0$ and $a(Y)=1$.
The base is $C\simeq\PP^1$. Every smooth fibre $F$ has
\[
 c_1(K_F)=0\text{ in }H^2(F;\Q),\qquad
 c_2(F)=0\text{ in }H^4(F;\Z),\qquad
 \chi(F,\OO_F)=0,
\]
and contains no exceptional curve of the first kind. A very general
smooth fibre has $\kappa(F)\leq0$. Its type belongs to the following
list: a complex two-torus, a bielliptic surface, a primary or secondary
Kodaira surface, a primary or secondary Hopf surface, or an Inoue
surface. Every listed type has Euler characteristic zero. This is a
preliminary necessary list. Statement C.3.7 below excludes the
bielliptic entry, and C.3.8 makes the very general Hopf fibre have
algebraic dimension zero. C.3.9 and C.3.10 exclude both Kodaira
entries. Statement C.3.6 realizes both primary and secondary Hopf
entries, and D.2.7 realizes a fixed abelian surface as general
fibre. Statement C.3.12 excludes the Inoue entry. Thus the
final very general list consists exactly of two-tori and
primary or secondary Hopf surfaces of algebraic dimension zero.
\end{proposition}
\begin{proof}
Connected fibres give $f_*\Q_Y=\Q_C$ and the degree-one edge injection
$H^1(C;\Q)\hookrightarrow H^1(Y;\Q)$. Therefore the genus of $C$ is
zero. The holomorphic tangent sequence on a smooth fibre, with the
last map induced by the differential of $f$, is
\[
0\longrightarrow T_F\longrightarrow T_Y|_F
 \xrightarrow{df} \OO_F\otimes T_{C,f(F)}\longrightarrow0.
\]
The normal factor is a trivial holomorphic line bundle. Thus
$K_F=K_Y|_F$ and $c_2(F)=c_2(Y)|_F$. Rational $H^2(Y)$ and $H^4(Y)$
vanish, and $H^4(F;\Z)=\Z$ has no torsion. These facts give the two
Chern-class assertions. Noether's formula gives $\chi(\OO_F)=0$.
An exceptional curve $E\simeq\PP^1$ with $E^2=-1$ would have
$K_F\cdot E=-1$ by adjunction, contradicting the rational vanishing of
$c_1(K_F)$.

For completeness we derive the Kodaira-dimension restriction directly.
Fix $m>0$ and put $G_m=f_*(K_Y^m)$. This is a coherent torsion-free
sheaf on the smooth curve $C$: a section killed by a nonzero base
function vanishes on a dense open subset of $Y$ and hence everywhere.
It is consequently a vector bundle. On the nonempty open set of
cohomological base change its fibre is $H^0(F,K_F^m)$. A rank at least
two would give a meromorphic map
\[
 Y\dashrightarrow \PP_C(G_m^*)
\]
whose restriction to a general fibre is defined by at least two
independent sections of a line bundle and has positive-dimensional
image. Its image has dimension at least two. The target is projective:
a vector bundle on a compact Riemann surface is algebraic, and its
projectivization is a projective variety. The image is therefore
projective and its rational functions pull back to a subfield of
$\mathcal M(Y)$ of transcendence degree at least two. This contradicts
$a(Y)=1$. Thus $\rk G_m\leq1$ for every $m$. Removing the countable
union of the exceptional base-change sets gives the claimed very
general-fibre assertion.

We now apply the Enriques--Kodaira classification in the form of
\cite[Theorem 3.1]{Peters1990}. The surface is minimal by the
adjunction argument. Rational and ruled surfaces have a rational
curve on which the canonical class has nonzero degree, so they are
excluded. K3 and Enriques surfaces have Euler numbers $24$ and $12$,
respectively. Properly elliptic and general-type surfaces have
positive Kodaira dimension. The remaining K\"ahler possibilities are
tori and bielliptic surfaces; the remaining Kodaira-dimension-zero
non-K\"ahler possibilities are the two types of Kodaira surfaces.
For class VII, $b_1(F)=1$, and $e(F)=0$ gives $b_2(F)=0$. The
classification of this case gives Hopf and Inoue surfaces; its
precise scope, including finite covers, is stated in
\cite[p.~2]{Teleman2008} and \cite[Theorem 1]{BBK2019}.
No assertion about torsion in $H_1(F;\Z)$ was used.
\end{proof}



\begin{proposition}[C.3.2; Refuted surface assertion, with explicit example]
\label{st:C.3.2}
The assertion that every secondary Kodaira surface has torsion in
$H_1(-;\Z)$ is false. There is a secondary Kodaira surface $Q$ with
\[
H_1(Q;\Z)=\Z,\qquad H^2(Q;\Z)=0,
\]
whose canonical bundle has holomorphic order six.
\end{proposition}
\begin{proof}
Let $H$ be the real Heisenberg group in coordinates $(x,y,z)$ with law
\[
(x,y,z)(x',y',z')
=\left(x+x',y+y',z+z'+\frac{xy'-yx'}2\right).
\]
Let $\Gamma\subset H$ be generated by
$a=(1,0,0)$, $b=(0,1,0)$ and $c=(0,0,1)$. Thus
\[
\Gamma=\{(x,y,z):x,y\in\Z,\ z-xy/2\in\Z\},
\qquad [a,b]=c.
\]
It is a discrete cocompact subgroup. On $H\times\mathbb R_t$ write
\[
\tau=e^{\pi i/3},\quad
\eta=dz-\frac{x\,dy-y\,dx}{2},\quad
\alpha=dx+\tau\,dy,\quad
\beta=\eta-\frac{dx+dy}{2}+i\,dt.
\]
These two complex one-forms and their conjugates are linearly
independent over $\C$. Declaring $\alpha,\beta$ to have type $(1,0)$
defines a left-invariant almost complex structure. It is integrable:
$d\alpha=0$ and $d\beta=-dx\wedge dy$ has type $(1,1)$, so the
ideal generated by the $(1,0)$ forms has no $(0,2)$ obstruction.
The structure descends to
$P=\Gamma\backslash H\times(\mathbb R/\Z)$.
The nonzero form $\alpha\wedge\beta$ is $\bar\partial$-closed and
trivializes $K_P$. The projection induced by $x+\tau y$ is a
holomorphic elliptic bundle over $\C/(\Z+\tau\Z)$, with nonzero
central Chern class $[a,b]=c$. Its first Betti number is three, so
$P$ is a primary Kodaira surface.

Define the Lie-group automorphism
\[
\sigma(x,y,z)=(-y,x+y,z-y/2).
\]
It preserves $\Gamma$: the difference between the new central
coordinate and half the product of the new horizontal coordinates
is
\[
(z-xy/2)+xy+y(y-1)/2.
\]
This is integral on $\Gamma$. The inverse preserves the lattice too,
and direct substitution gives $\sigma^6=1$. Moreover
$\sigma^*\alpha=\tau\alpha$ and $\sigma^*\beta=\beta$.
Consequently
\[
g:P\longrightarrow P,\qquad
g([x,y,z],t)=([\sigma(x,y,z)],t+1/6)
\]
is a holomorphic automorphism of order six. Every nonidentity power
is fixed-point-free because of its translation of the $t$-circle.
The quotient $Q=P/\langle g\rangle$ is therefore a compact complex
surface and $P\to Q$ is finite \emph{\'etale} of degree six.

The contractible manifold $H\times\mathbb R$ is its universal cover.
The deck group is $\Gamma\rtimes_\sigma\Z$, where the infinite
generator is the lift $(\sigma,t\mapsto t+1/6)$; its sixth power is
the original integral translation in $t$. Since
$\Gamma^{\rm ab}=\Z^2$ and the induced matrix is
\[
A=\begin{pmatrix}0&-1\\1&1\end{pmatrix},\qquad
\det(A-I)=1,
\]
abelianization gives $H_1(Q;\Z)=\Z$.
The finite cover and $e(P)=0$ give $e(Q)=0$. Poincar\'e duality then
gives $b_2(Q)=0$, and the universal coefficient theorem gives
$H^2(Q;\Z)=0$ because $H_1(Q;\Z)$ has no torsion.
Finally $g^*(\alpha\wedge\beta)=\tau\alpha\wedge\beta$.
Every holomorphic section of $K_P^m$ is a constant multiple of
$(\alpha\wedge\beta)^m$, by compactness. Such a section descends
precisely when $6$ divides $m$. Thus $K_Q$ has holomorphic order six.
The surface has $b_1=1$ and a primary Kodaira surface as a finite
\'etale cover, so it is secondary Kodaira. The matrix and lattice
calculations have been checked exactly by the attached program.
\end{proof}

\begin{proposition}[C.3.3; Refuted surface assertion, with explicit example]
\label{st:C.3.3}
The assertion that a secondary Hopf surface has torsion in
$H_1(-;\Z)$ is false. There is a secondary Hopf surface $Q_H$ with
$H_1(Q_H;\Z)=\Z$ and $H^2(Q_H;\Z)=0$.
\end{proposition}
\begin{proof}
Write a quaternion as $(q_0,q_1,q_2,q_3)$ and put
$\varphi=(1+\sqrt5)/2$. Let $G$ be the following set of 120 unit
quaternions: the eight signed coordinate units; the sixteen points
$(\pm1,\pm1,\pm1,\pm1)/2$; and the 96 points obtained by even
coordinate permutations of
$(0,\pm1,\pm\varphi,\pm\varphi^{-1})/2$, with independent signs.
Quaternion multiplication makes this the binary icosahedral group.
For an entirely finite verification, the attached program enumerates
these points in $\Q(\sqrt5)^4$, verifies unit norm and all 14,400
multiplication-table entries, and computes its commutator subgroup.
Every element is a commutator, in particular $G^{\rm ab}=0$.
This is an exact computation over rational pairs representing
$\Q(\sqrt5)$.

Use the faithful holomorphic representation
\[
(a,b,c,d)\longmapsto
\begin{pmatrix}a+ib&c+id\\-c+id&a-ib\end{pmatrix}\in SU(2)
\]
and let the generator of $\Z$ act on $\C^2\setminus\{0\}$ by
$v\mapsto2v$. The action of $\Z\times G$ is free. An equality
$2^ngv=v$ first gives $n=0$ by norms, and a unit quaternion whose
matrix has eigenvalue one is the identity. Proper discontinuity
follows from the radial scaling and finiteness of $G$.
Polar coordinates identify the quotient smoothly with
\[
Q_H=(\C^2\setminus\{0\})/(2^{\Z}\times G)
\simeq (\mathbb R/(\log2)\Z)\times(S^3/G).
\]
It is compact, and its universal cover is $\C^2\setminus\{0\}$.
The fundamental group is $\Z\times G$, whose abelianization is
$\Z$. The quotient by $2^{\Z}$ alone is a primary Hopf surface
covering $Q_H$ with degree 120. The fundamental group of $Q_H$ is
noncyclic, so $Q_H$ is secondary Hopf. The product description gives
$e(Q_H)=0$, and the same duality and universal-coefficient argument
as before gives $H^2(Q_H;\Z)=0$.
\end{proof}

These two explicit finite-cover maps locate two additional failures
in the surface-level exclusions used in the proof of CDP20,
Lemma 4.3. They do not construct these surfaces as fibres of a
threefold satisfying the standing hypotheses. Such a realization or
a separate exclusion remains part of C.3.

\begin{proposition}[C.3.4; Proved]\label{st:C.3.4}
Let $Y$ be a compact complex threefold with $b_2(Y)=0$, and let
$f:Y\to C$ be proper, holomorphic, surjective, with connected fibres.
Every irreducible component of a reducible reduced fibre has no
compact K\"ahler resolution. In particular, a fibre all of whose
components admit compact K\"ahler resolutions is irreducible.
\end{proposition}
\begin{proof}
Let $D$ be a component of a reducible fibre. The incidence graph of
the components is connected, so some other component $D'$ meets $D$.
Both are Cartier divisors on the smooth threefold. Let
$\mu:\widehat D\to D\hookrightarrow Y$ be a resolution by a compact
smooth surface. Pullback of the canonical section of $\OO_Y(D')$
is a nonzero holomorphic section of $\mu^*\OO_Y(D')$. Its zero
divisor $E$ is effective and nonzero: the intersection $D\cap D'$
is a nonempty divisor on the irreducible surface $D$, and pullback
to the resolution preserves its nonzero vanishing order along a
component dominating that intersection. On the other hand
$c_1(\mu^*\OO_Y(D'))=0$ in $H^2(\widehat D;\R)$.
A K\"ahler form $\omega$ on $\widehat D$ would give
\[
0=\int_{\widehat D}c_1(\OO_{\widehat D}(E))\wedge\omega
 =\int_E\omega>0,
\]
where positive multiplicities are retained in the last integral.
This contradiction proves the assertion for every component.
\end{proof}

\begin{corollary}[C.3.5; Proved]\label{st:C.3.5}
For every threefold and fibration in C.3.1, including every type on
its necessary list, the Euler-number identity of Lead 5 is
\[
\int_Y c_3(Y)=e(Y)=\sum_{c\in\Delta}e\bigl((Y_c)_{\rm red}\bigr),
\]
where $\Delta$ contains the finitely many critical values.
The identity alone assigns no sign to the summands.
\end{corollary}
\begin{proof}
All the listed smooth surfaces have Euler number zero by C.3.1.
Additivity of compactly supported Euler characteristic for the
stratification by the smooth-fibre locus and the finitely many
exceptional fibres gives
$e(Y)=e_c(C\setminus\Delta)e(F)+\sum_{c\in\Delta}e(Y_c)$.
The first term vanishes. A scheme-theoretic fibre and its reduction
have the same underlying topological space, so their topological
Euler numbers agree. Chern--Gauss--Bonnet gives the first equality.
\end{proof}

\begin{example}[C.3.6; Proved actual occurrence of both Hopf types]
\label{st:C.3.6}
There are compact complex threefolds with $b_1=b_2=0$ and $a=1$
whose holomorphic algebraic reduction is a smooth bundle over
$\PP^1$ with primary Hopf fibre. There are also such threefolds
with secondary Hopf fibre. In these examples every fibre has Euler
number zero and $c_3=0$.
\end{example}
\begin{proof}
Put $a_1=1$, $a_2=\sqrt2$ and
$W=(\C^2\setminus\{0\})_z\times(\C^2\setminus\{0\})_w$.
The additive complex Lie group $\C$ acts holomorphically by
\[
t\cdot(z,w_1,w_2)
 =\bigl(e^t z,e^{ia_1t}w_1,e^{ia_2t}w_2\bigr).
\tag{C.H.1}
\]
This action is free. Indeed $e^t z=z$ implies $t=2\pi in$ with
$n\in\Z$; a nonzero $w_j$ then gives $e^{-2\pi a_jn}=1$ and
therefore $n=0$. It is proper as well. If a point and its translate
range in fixed compact subsets of $W$, the $z$-norm bounds
$\operatorname{Re}t$. Writing $s=\operatorname{Im}t$, the identity
\[
\|t\cdot w\|^2
 =e^{-2a_1s}|w_1|^2+e^{-2a_2s}|w_2|^2
\]
and the positive constants $a_1,a_2$ bound $s$ both above and
below, since both $w$-norms are bounded away from zero. The inverse
image of a compact subset under the action graph is consequently
closed and bounded in the group variable and compact in the other
variable. This proves properness.

Every orbit meets the real submanifold
$\{\|z\|=\|w\|=1\}=S^3\times S^3$ exactly once.
The first equation fixes $\operatorname{Re}t=-\log\|z\|$.
The displayed squared $w$-norm is strictly decreasing in $s$,
with limits $+\infty$ and $0$, so the second equation has one
solution. Its derivative is strictly negative, and the implicit
function theorem makes the resulting cross-section smooth.
Thus the complex quotient $Y=W/\C$ is a smooth compact complex
threefold diffeomorphic to $S^3\times S^3$. In particular
$b_1(Y)=b_2(Y)=0$ and $e(Y)=0$.

The map
\[
f:Y\longrightarrow\PP^1,\qquad [z,w]\longmapsto[z_1:z_2]
\]
is well-defined and holomorphic. Choose a holomorphic nonzero
frame of the tautological line on a small base disc. Scaling the
$z$-coordinate to this frame leaves exactly the transformations
$t=2\pi in$, which act on $w$ by powers of
\[
Q=\operatorname{diag}(q_1,q_2),\qquad
q_j=e^{-2\pi a_j}.
\]
This identifies $f$ over that disc with the product having fibre
$H=(\C^2\setminus\{0\})/\langle Q\rangle$.
The quotient is a primary Hopf surface, and these local products
prove that $f$ is a proper holomorphic submersion with connected
fibres.

We verify $a(H)=0$ directly. A meromorphic function on $H$ lifts
to a $Q$-invariant meromorphic function on $\C^2\setminus\{0\}$.
It extends meromorphically across the origin: the closure of its
polar divisor is analytic by the codimension-two extension theorem
for divisors, is locally principal on the smooth ball, and
multiplying by a local defining function reduces extension to
Hartogs' theorem for a holomorphic function. Write the germ as
$g/h$ with relatively prime convergent power series. The identity
$(g\circ Q)h=g(h\circ Q)$ gives
\[
g\circ Q=ug,\qquad h\circ Q=uh
\]
for one holomorphic unit $u$: divisibility follows from relative
primality, and preservation of the vanishing order makes the
quotient a unit. Write $u/u(0)=e^\ell$ with $\ell(0)=0$ and set
\[
v(w)=\exp\left(\sum_{n\ge0}\ell(Q^nw)\right).
\]
The sum converges normally on a sufficiently small ball, because
$Q$ is a strict contraction. It gives
$v(Qw)/v(w)=u(0)/u(w)$. Consequently both $vg$ and $vh$ are
eigenfunctions of pullback by $Q$ with the same eigenvalue $u(0)$.
The monomial $w_1^rw_2^s$ has eigenvalue
$e^{-2\pi(r+s\sqrt2)}$. These numbers are distinct for distinct
$(r,s)\in\Z_{\ge0}^2$. Comparing power-series coefficients shows
that each nonzero eigenfunction is a single monomial and that the
two eigenfunctions are proportional. The meromorphic function is
therefore constant, first near the origin and then everywhere by
analytic continuation.

Every meromorphic function on $Y$ is constant on each general
fibre of $f$. On a product chart, restricting to a local section
gives the descended meromorphic function on the base; equality on
the fibres and the identity theorem identify its pullback with
the original function. These local functions agree on overlaps
and extend at the finitely many possible poles on the base.
Thus $\mathcal M(Y)=f^*\mathcal M(\PP^1)$, proving $a(Y)=1$
and proving that $f$ is the holomorphic algebraic reduction.

For any integer $m>1$, let $\mu_m$ act on $W$ by
$(z,w)\mapsto(z,\zeta w)$. It commutes with (C.H.1) and induces
a free action on $Y$. A fixed orbit would give $t=2\pi in$ from
the $z$-coordinate; equality of norms in a nonzero $w_j$ then
gives $n=0$ and $\zeta=1$. The quotient
$Y_m=Y/\mu_m$ is a smooth compact complex threefold. Transfer
injects $H^j(Y_m;\Q)$ into $H^j(Y;\Q)$, so $b_1(Y_m)=b_2(Y_m)=0$.
The simply connected cover $Y\simeq S^3\times S^3$ gives
$\pi_1(Y_m)=\mu_m$, and the universal coefficient theorem gives
$H^2(Y_m;\Z)\simeq\Z/m\Z$: its free part is zero by the
rational calculation, and its torsion is
$\operatorname{Ext}(H_1(Y_m;\Z),\Z)$. Thus these secondary
examples meet exactly the rational hypotheses of C.3; the primary
example also has $H^1(Y;\Z)=H^2(Y;\Z)=0$.
The induced map to $\PP^1$ has fibre
\[
H_m=(\C^2\setminus\{0\})/
       \bigl(\langle Q\rangle\times\mu_m\bigr).
\]
Its universal cover is $\C^2\setminus\{0\}$ and its fundamental
group is $\Z\times\mu_m$, which is noncyclic. It is therefore a
secondary Hopf surface, with the displayed primary cover of degree
$m$. Any meromorphic function on $H_m$ pulls back to one on $H$,
so $a(H_m)=0$. The same local descent proves
$\mathcal M(Y_m)=f_m^*\mathcal M(\PP^1)$ and $a(Y_m)=1$.
Finally transfer of Euler characteristic under the finite covering,
or the smooth-fibration Euler formula, gives $e(Y_m)=0$.
\end{proof}

\begin{proposition}[C.3.7; Proved exclusion of bielliptic fibres]
\label{st:C.3.7}
A compact complex threefold $Y$ with $b_1(Y)=b_2(Y)=0$ has no
proper connected-fibre holomorphic map to a compact curve whose
general fibre is bielliptic.
\end{proposition}
\begin{proof}
We use precisely the relative Albanese construction of
\cite[Proposition~2, pp.~605--608]{Kawai1969}, and not its subsequent
corollary. Applied to bielliptic fibres, which have
$p_g=0$ and $h^0(\Omega^1)=1$, it gives a compact complex surface
$B$, a holomorphic map $h:B\to C$, and a meromorphic map
$\Phi:Y\dashrightarrow B$ over $C$. On the complement $C^\circ$
of finitely many values, $B^\circ\to C^\circ$ is a smooth family
of elliptic curves and
\[
 \Phi|_{F_c}:F_c\longrightarrow B_c
\]
is the Albanese morphism of the bielliptic surface, with connected
elliptic fibres. To specify the comparison being imported, the
construction forms the relative period-lattice quotient, glues its
local Albanese maps by fibre translations, and compactifies at each
missing value by a normalized finite base change and the projective
space associated with the direct images of a divisor which is the
inverse image of a local Albanese section. In making this last
step equivariant, use the sum of the finite orbit of that divisor;
a selected section need not itself be invariant under the covering
group. Resolve the normalized base change first if needed, so the
divisor defines a line bundle. The orbit sum restricts on every
general fibre to the pullback of a positive-degree divisor on its
Albanese elliptic curve; sufficiently large powers give the same
Albanese image. Its naturally linearized line bundle and direct
images carry the finite covering-group action, and their quotient
glues the construction to $B^\circ$. These are the maps asserted by
Proposition~2, with that equivariant choice made explicit. The
later assertion in Kawai's corollary about multiple fibres is not
used here.

Resolve $B$ over $C\setminus C^\circ$ and resolve the induced
meromorphic map by a sequence of blowups with smooth centres
supported over those same values. This produces a diagram
\[
 \widetilde Y\xrightarrow{\ g\ }\widetilde B,
 \qquad \pi:\widetilde Y\longrightarrow Y,
 \qquad (h\circ\rho)\circ g=f\circ\pi,
\]
where $\rho:\widetilde B\to B$ is a resolution and both spaces
are smooth and compact. The map $g$ has connected fibres:
its general fibres are the connected Albanese fibres just described,
and its Stein factorization has finite degree one over the normal
surface $\widetilde B$. Hence
$H^1(\widetilde B;\Q)\hookrightarrow H^1(\widetilde Y;\Q)$.
Blowing up smooth centres does not change first cohomology, so
$b_1(\widetilde B)=0$. The compact-surface K\"ahler criterion
\cite[Theorem~1.1]{Peters1990} therefore gives a K\"ahler form
$\omega_B$ on $\widetilde B$.

The blowup cohomology formula gives
\[
 H^2(\widetilde Y;\R)
 =\pi^*H^2(Y;\R)\oplus\bigoplus_j\R[E_j]
 =\bigoplus_j\R[E_j],
\]
where the $E_j$ are the exceptional prime-divisor classes (with
the corresponding total transforms for successive blowups).
Every $E_j$ is supported above one of the removed values in $C$.
Consequently its restriction to a general smooth fibre $F_c$,
identified with its unchanged inverse image in $\widetilde Y$,
is zero. In particular
\[
 [g^*\omega_B]|_{F_c}=0\quad\hbox{in }H^2(F_c;\R).
\]
Choose a K\"ahler form $\omega_F$ on the bielliptic surface $F_c$.
The Albanese morphism $g|_{F_c}$ is a submersion to the elliptic
curve $\widetilde B_c$. Thus its pullback of $\omega_B$ is a
nonzero semipositive $(1,1)$-form and
\[
 \int_{F_c}g^*\omega_B\wedge\omega_F>0.
\]
The vanishing cohomology class makes the same integral zero by
Stokes' theorem, a contradiction. All maps in this argument are
the displayed relative Albanese map and its modifications; no
comparison between the two elliptic fibrations of a bielliptic
surface is needed.
\end{proof}

\begin{proposition}[C.3.8; Proved exclusion of elliptic Hopf general fibres]
\label{st:C.3.8}
The very general Hopf fibre of the algebraic reduction in C.3.1
has algebraic dimension zero. In particular neither primary nor
secondary Hopf surfaces of algebraic dimension one can be its
general fibre. This exclusion uses no restriction-surjectivity
assertion from CDP98 or CDP20.
\end{proposition}
\begin{proof}
We first verify that a Hopf surface $H$ with $a(H)=1$ satisfies
$\kappa(H,K_H^{-1})=1$. On a primary Hopf cover the contraction
normal forms are diagonal, or
\[
A(z,w)=(\alpha z+\lambda w^n,\beta w),\qquad
\lambda\ne0,\quad \alpha=\beta^n.
\]
The invariant bivectors $zw\,\partial_z\wedge\partial_w$ in the
diagonal case and $w^{n+1}\,\partial_z\wedge\partial_w$ in the
second case give a nonzero effective anticanonical divisor.
The invariance is an equality under the derivative of the
displayed contraction: its determinant is $\alpha\beta$, and
the displayed coefficient transforms by the same scalar.
These normal forms and the elliptic-fibration description are
also used in \cite[Proposition~3.1 and its proof]{SS2002}.

The algebraic reduction of an elliptic Hopf surface is a
holomorphic map to $\PP^1$, with smooth reductions of all fibres.
Every curve is vertical: its intersection with a general fibre
is zero since $H^2(H;\R)=0$, whereas a horizontal curve would
have positive intersection equal to its degree over the base.
Thus the nonzero anticanonical divisor is a positive sum of
reduced fibres. If $m_iD_i$ is the full fibre supported on
$D_i$, each $\OO(m_iD_i)$ is the pullback of $\OO_{\PP^1}(1)$.
A common multiple of the $m_i$ therefore makes a positive power
of $K_H^{-1}$ the pullback of a positive-degree line bundle on
$\PP^1$. Its Iitaka dimension is one.

For a secondary Hopf surface take a normal finite primary cover
$p:H'\to H$. The cover is unramified, so
$K_{H'}^{-1}=p^*K_H^{-1}$, and $a(H')=a(H)=1$ by invariance of
transcendence degree under a finite map. The anticanonical
section ring of $H$ is the invariant subring of that of $H'$.
The latter is integral over the invariant subring: a homogeneous
section $s$ satisfies the monic orbit polynomial
$\prod_{g\in\operatorname{Gal}(H'/H)}(T-gs)$. Their fraction
fields have the same transcendence degree, proving
$\kappa(H,K_H^{-1})=1$ also in this case.

Now, for each positive integer $m$, the locus
\[
Z_m=\{c\in C^\circ:
 h^0(Y_c,K_Y^{-m}|_{Y_c})\ge2\}
\]
is a closed analytic subset of the smooth base $C^\circ$.
If the very general Hopf fibre had algebraic dimension one,
the preceding calculation would make the union of these
countably many loci contain the complement of a countable
subset of $C^\circ$. A proper closed analytic subset of a
curve is discrete and countable, so some $Z_m$ is the whole
connected smooth base. Coherent base change then gives
$\operatorname{rank}f_*K_Y^{-m}\ge2$.
The relative linear system defines a meromorphic map of $Y$
to $\PP(f_*K_Y^{-m})$ over $C$ with positive-dimensional image
on a general fibre. Its image has dimension at least two.
The projective bundle over the compact curve is projective,
so pulling back its meromorphic functions gives $a(Y)\ge2$,
contrary to the hypothesis. In fact the same argument shows
that, when the smooth fibres are Hopf surfaces, the values with
$a(Y_c)=1$ form at most a countable set.
\end{proof}

\begin{theorem}[C.3.9; Proved exclusion of secondary Kodaira general fibres]
\label{st:C.3.9}
A compact complex threefold with $b_1(Y)=b_2(Y)=0$ and
$a(Y)=1$ cannot have a holomorphic algebraic reduction
$f:Y\to\PP^1$ whose general fibre is a secondary Kodaira
surface. Rational Betti-number vanishing suffices; no integral
torsion-freeness of the fibre is assumed.
\end{theorem}
\begin{proof}
Suppose that the smooth general fibres $S_c$ are secondary
Kodaira surfaces. They have $h^1(\OO_{S_c})=b_1(S_c)=1$.
The nonzero integral transgression section of D.1.9 has
coefficient image $\beta\in H^1(Y,\OO_Y)$ and gives a
nowhere-zero section of $\mathscr E=R^1f_*\OO_Y$.
This coherent sheaf is a line bundle by E4. Thus
$\mathscr E=\OO_{\PP^1}$ and
$H^1(Y,\OO_Y)=\C\beta$. For every $z\in\C$ set
\[
 M_z=\exp(2\pi\mathrm i z\beta)\in\Pic^0(Y).
\tag{C.SK.1}
\]
Its restriction satisfies $M_{z+1}|_{S_c}\simeq M_z|_{S_c}$,
because $\beta|_{S_c}$ is the coefficient image of an
integral class. It is nonzero as a tangent vector of
$\Pic^0(S_c)$ by D.1.7. The exponential sequence therefore
identifies $\Pic^0(S_c)$ with $\C^*$ and makes (C.SK.1), on
each smooth fibre, a nonconstant homomorphism from
$\C/\Z$ to this group.

We specify the primary cover and the invariant used below.
Let $m>1$ be the order of $K_{S_c}$, and let
\[
 p_c:P_c\longrightarrow S_c
\]
be its connected cyclic canonical cover, with deck group
$G=\mu_m$. It is finite \emph{\'etale}. Its canonical bundle
is trivial, and it is non-K\"ahler, because a finite
\emph{\'etale} quotient of a K\"ahler surface is K\"ahler
by averaging a K\"ahler form. The classification of
surfaces with trivial canonical bundle consequently makes
$P_c$ a primary Kodaira surface. Write
$\pi_c:P_c\to E_c$ for its Albanese elliptic bundle, and
$d>0$ for its topological degree. These covers exist in
holomorphic families locally on the smooth base: the
evaluation of a local frame of $f_*K_{Y/\PP^1}^{m}$
trivializes $K_{Y/\PP^1}^{m}$, since it is nowhere zero on
each fibre, and taking its cyclic root gives the required
finite \emph{\'etale} covering family. The integers $m,d$
are locally constant and hence constant on the connected
smooth base.

The exact primary-surface invariant we use is
\[
 Q_c=\Pic^0(P_c)/\pi_c^*\Pic^0(E_c)\simeq\C^*.
\tag{C.SK.2}
\]
The Picard pullback in the denominator is injective:
$\pi_{c*}\OO_{P_c}=\OO_{E_c}$ and the projection formula
show that a line with trivial pullback is trivial. For
any line $L_c$ of degree $d$ on $E_c$, the pullback
$\pi_c^*L_c$ belongs to $\Pic^0(P_c)$, and its class in
$Q_c$ is independent of $L_c$. Choose the group
isomorphism $q_c:Q_c\to\C^*$ that places this
positive-degree class in the punctured unit disc, and
denote its value by $\alpha(c)$.
This fixes the possible inversion ambiguity of the
group coordinate. The descriptions in
\cite[Theorem~4 and Remark~1]{Borcea1984} and
\cite[\S1, pp.~2--3]{SS2002} identify $\alpha$ as the
intrinsic holomorphic punctured-disc coordinate of the
moduli of primary Kodaira surfaces. More explicitly,
$P_c$ is the associated elliptic bundle of a line of
degree $d$ with fibre $\C^*/\langle\alpha(c)\rangle$;
quotienting by $\pi_c^*\Pic^0(E_c)$ removes the choice
of that line within its degree. In a holomorphic
family the local integral marking and the coefficient
map give this quotient a holomorphic $\C^*$ coordinate,
and the chosen positive-degree orientation is locally
constant. This is also immediate in Borcea's complete
period family. Thus $\alpha(c)$ is holomorphic.

The locally constructed canonical covers have the same
isomorphism class on overlaps, so this intrinsic
function glues on the smooth base. It has
$0<|\alpha(c)|<1$. It extends across the finitely many
omitted values by the removable singularity theorem
and is constant on $\PP^1$. Write its nonzero constant
value as $\alpha$.

Consider the exact homomorphism of connected Picard
groups
\[
 \Pic^0(S_c)\xrightarrow{\ p_c^*\ }\Pic^0(P_c)
          \longrightarrow Q_c\xrightarrow{\ q_c\ }\C^*.
\tag{C.SK.3}
\]
Its differential is nonzero. Indeed finite-cover trace
identifies $H^1(S_c,\OO_{S_c})$ with the $G$-invariant
line in $H^1(P_c,\OO_{P_c})$. The other line,
$\pi_c^*H^1(E_c,\OO_{E_c})$, has nontrivial $G$ action:
otherwise the unique holomorphic one-form of $P_c$
would descend to $S_c$, whereas a secondary Kodaira
surface has no holomorphic one-form. Hence the invariant
line has zero intersection with the denominator's
tangent line in (C.SK.2). This proves the claim.
Consequently (C.SK.3) is a nonzero isogeny of $\C^*$
groups. Composing it with (C.SK.1) gives
\[
 q_c\bigl([p_c^*(M_z|_{S_c})]\bigr)
       =\exp(2\pi\mathrm i k z)
       \qquad(k\in\Z\setminus\{0\}).
\tag{C.SK.4}
\]
Here integrality follows from the period $z\mapsto z+1$;
every holomorphic group homomorphism $\C^*\to\C^*$ is
an integral power. The integer $k$ is locally constant
in the holomorphic families just described, and is
therefore a single integer on the connected smooth base.

Let $\rho_c:S_c\to\PP^1$ be its elliptic algebraic
reduction. The induced finite map $r_c:E_c\to\PP^1$
satisfies $\rho_c p_c=r_c\pi_c$. Write $h=\deg r_c$;
this positive integer is constant. Choose one positive
integer $N$ killing the torsion group $H^2(S_c;\Z)$
and satisfying $d\mid hN$. This choice is uniform
because the smooth fibres are differentiably locally
trivial. Then
$A_c=\rho_c^*\OO_{\PP^1}(N)$ belongs to $\Pic^0(S_c)$,
and (C.SK.2) gives
\[
 q_c([p_c^*A_c])=\alpha^{hN/d}.
\tag{C.SK.5}
\]
This identity is fibrewise and uses only degrees; it
does not require the divisors $\rho_c^{-1}(x)$ to
extend across a singular fibre of $f$.

Choose $z_0\in\C$ with
$\exp(2\pi\mathrm i k z_0)=\alpha^{hN/d}$ and put
$M=M_{z_0}$. By (C.SK.4)--(C.SK.5), the difference of
$p_c^*(M|_{S_c})$ and $p_c^*A_c$ is
$\pi_c^*R_c$ for some $R_c\in\Pic^0(E_c)$.
Both of the first two lines come from $S_c$ and are
$G$-invariant. Injectivity of $\pi_c^*$ implies that
$R_c$ is fixed by the induced nontrivial linear action
on $\Pic^0(E_c)$. An elliptic-curve linear automorphism
of finite order other than one has order
$r\in\{2,3,4,6\}$. If it fixes a point $R_c$, then
\[
 rR_c=(1+\zeta+\cdots+\zeta^{r-1})R_c=0
\]
in the elliptic group. Thus $R_c^{\otimes12}$ is
trivial. The norm of the finite cover satisfies
$\operatorname{Nm}_{p_c}(p_c^*L)=L^{\otimes m}$.
Applying it to the difference now proves the exact
line-bundle identity
\[
 M^{\otimes12m}|_{S_c}
       \simeq \rho_c^*\OO_{\PP^1}(12mN)
\tag{C.SK.6}
\]
for every smooth general fibre.

The right side has at least two linearly independent
sections, by pullback from $\PP^1$. Coherent generic
base change gives
$\operatorname{rank}f_*M^{\otimes12m}\ge2$.
Its relative linear system has a positive-dimensional
image on a general fibre and therefore a surface
image in the projective bundle over $\PP^1$.
Pullback of meromorphic functions from this image
gives $a(Y)\ge2$, the desired contradiction.
This proves the exclusion without the arbitrary-twist
cohomological-flatness assertion used in the printed
Kodaira exclusion.
\end{proof}

\begin{theorem}[C.3.10; Proved exclusion of primary Kodaira general fibres]
\label{st:C.3.10}
A compact complex threefold with $b_1(Y)=b_2(Y)=0$ and
$a(Y)=1$ cannot have a holomorphic algebraic reduction
$f:Y\to\PP^1$ whose general fibre is a primary Kodaira
surface. Together with C.3.9 this excludes both Kodaira
types using only rational Betti-number vanishing.
\end{theorem}
\begin{proof}
Let $B=\PP^1\setminus\Sigma$ be the smooth base and write
$P_c=f^{-1}(c)$ and $\pi_c:P_c\to E_c$ for the canonical
elliptic algebraic reduction of a smooth fibre. The degree
$d>0$ is constant. All the smooth fibres have the fixed
primary Kodaira differentiable type; every complex structure
on that type is again a primary Kodaira structure
\cite[pp.~373--374]{Borcea1984}.
Use the intrinsic quotient and its orientation from (C.SK.2):
\[
 Q_c=\Pic^0(P_c)/\pi_c^*\Pic^0(E_c),\qquad
 q_c:Q_c\xrightarrow{\simeq}\C^*.
\]
The class of the pullback of a degree-$d$ line has coordinate
$\alpha(c)\in\Delta^*$. The holomorphic moduli comparison
in the proof of C.3.9 applies directly here, so this bounded
function extends to $\PP^1$ and has a single nonzero
constant value $\alpha$.

Put $\mathscr E=R^1f_*\OO_Y$ and
$\mathscr G=R^2f_*\OO_Y$. The primary Kodaira case of
D.1.12 proves, on the original threefold and every full fibre,
\[
 \operatorname{rank}\mathscr E=2,\qquad
 \operatorname{rank}\mathscr G=1,\qquad
 \bigwedge^2\mathscr E\xrightarrow{\simeq}\mathscr G.
\tag{C.PK.1}
\]
Both sheaves are locally free and satisfy base change.
Let $\gamma$ be the integral transgression section of D.1.9,
and let $\beta\in H^0(\PP^1,\mathscr E)=H^1(Y,\OO_Y)$
be its coefficient image. This section is nowhere zero.
For $u\in H^1(P_c,\OO_{P_c})$ define the complex-linear
functional $\lambda_c$ by the exact identity
\[
 q_c\bigl([\exp(2\pi\mathrm i z u)]\bigr)
     =\exp(2\pi\mathrm i z\lambda_c(u))
     \quad(z\in\C).
\tag{C.PK.2}
\]
Its kernel is the line
$A_c=\pi_c^*H^1(E_c,\OO_{E_c})$.
The function $k(c)=\lambda_c(\beta_c)$ takes integral
values, because $\beta_c$ is the coefficient image of
an integral class. It is holomorphic, by the actual relative
Picard quotient in (C.SK.2), hence is a constant integer $k$.

When $k\ne0$, choose $n>0$ with $nd\ge2$ and choose
$z\in\C$ with $\exp(2\pi\mathrm i kz)=\alpha^n$.
The ambient line $M=\exp(2\pi\mathrm i z\beta)$ then
has, on every smooth fibre, the same $Q_c$ class as the
pullback of a degree-$nd$ line on $E_c$. Its difference
from that pullback belongs to $\pi_c^*\Pic^0(E_c)$.
Thus
$M|_{P_c}=\pi_c^*L_c$ for a line $L_c$ of degree $nd$,
and $h^0(P_c,M|_{P_c})=nd\ge2$. The relative linear
system gives $a(Y)\ge2$, as at the end of C.3.9.
It remains to obtain the same contradiction when $k=0$.

In that case $\beta_c\in A_c$. The coefficient map
$H^1(E_c;\R)\to H^1(E_c,\OO_{E_c})$ is an isomorphism
of real vector spaces. Its pullback, together with the real
injectivity of D.1.7 on $P_c$, proves the exact implication
\[
 \beta_c\in A_c
 \quad\Longrightarrow\quad
 \gamma_c\in\pi_c^*H^1(E_c;\R).
\tag{C.PK.3}
\]
The lattice $\pi_c^*H^1(E_c;\Z)$ is saturated in
$H^1(P_c;\Z)$: the elliptic bundle induces a surjection
on fundamental groups, so the cokernel of the induced
injection on integral $H^1$ is torsion free. Thus $\gamma$
lies in this rank-two integral local system. Divide it by
its fixed positive divisibility to obtain a primitive fixed
section $\gamma_0$. On a simply connected open set choose
a complementary integral section $\delta$ with oriented
intersection number one. The ratio of their nonzero
coefficient images is a holomorphic elliptic period
\[
 \tau=\frac{\iota(\delta)}{\iota(\gamma_0)},
 \qquad \operatorname{Im}\tau>0,
\]
after choosing the sign convention of the ordered pair.
Monodromy fixes $\gamma_0$ and replaces $\delta$ by
$\delta+N\gamma_0$, $N\in\Z$. Therefore
$\exp(2\pi\mathrm i\tau)$ is a single-valued bounded
holomorphic function on $B$. It extends to $\PP^1$ and
is constant. Its logarithmic derivative gives $d\tau=0$,
and a nonzero integral translation cannot fix this constant
period. Consequently the Albanese integral $H^1$ local
system has identity monodromy, and $j(E_c)$ is constant.

Both moduli $(j,\alpha)$ and the degree $d$ are now
constant. Borcea's classification, in the precise form
\cite[\S1]{SS2002}, makes all smooth fibres biholomorphic
to one primary Kodaira surface $P$. Here is the actual
local product comparison we use. Borcea's complete
holomorphic marked period family \cite[Remark~1]{Borcea1984}
expresses our family, near any smooth base point, as the
pullback by a holomorphic map into his period domain
$\mathcal D\subset\PP(H^2(P;\C))$. By his Theorem~3,
constant isomorphism class means that the image of this
map lies in one orbit of the integral cohomology action.
That orbit is countable, since the action factors through
$\operatorname{GL}(H^2(P;\Z)/\operatorname{torsion})$.
A connected subset of a countable subset of a complex
manifold is one point: a real coordinate separating two
points would have a nontrivial interval as its connected
image. Therefore the period parameter is constant and
its pullback is a product. This proves the needed
holomorphic local triviality from the specified complete
family.

We compute its entire action on $H^2(P;\C)$, including
vertical shears. Write the fundamental group as
\[
 \Gamma=\langle x,y,w,u\mid
       w,u\text{ central},\ [x,y]=w^d\rangle.
\tag{C.PK.4}
\]
Its centre $\langle w,u\rangle$ is the fundamental group
of the elliptic fibre, and its quotient is the fundamental
group of the elliptic base. A biholomorphism inducing the
identity on the base $H^1$ induces the identity on this
quotient. The relation in (C.PK.4) then forces it to fix
$w$. On an elliptic fibre its linear part is a holomorphic
group automorphism. Such an automorphism acts by a complex
scalar on the universal cover; fixing the nonzero period
$w$ forces that scalar to be one. It therefore fixes both
$w$ and $u$. The remaining possible actions on $x,y$ are
exactly multiplication by central elements.

For clarity, their induced cohomology action can be read
off without discarding any of those elements. In the real
nilmanifold presentation choose invariant one-forms
$a,b,c,e$ with
\[
 da=db=de=0,\qquad dc=d\,a\wedge b.
\]
The four classes $[a\wedge c],[b\wedge c],
[a\wedge e],[b\wedge e]$ form a basis of
$H^2(P;\R)$, as in \cite[\S3]{Borcea1984}.
An automorphism of (C.PK.4) with the stated action extends
to the real nilpotent group: in logarithmic coordinates it
fixes the two central generators and adds central vectors
to the two base generators. Its pullback fixes $a,b$ and
replaces each of $c,e$ by itself plus a real linear
combination of $a,b$. Each of the four displayed classes
therefore changes only by a multiple of the exact form
$a\wedge b$. Since the universal cover of $P$ is
contractible, a map of $P$ and this group-model map with
the same induced fundamental-group map are homotopic.
This computes the action of the original biholomorphism
and proves it is the identity on $H^2$.
In particular every monodromy of our isotrivial family,
whose action on the base $H^1$ was just proved to be the
identity, acts identically on $H^2(P;\C)$.

We next construct a global section of $\mathscr G$ by an
actual morphism of complexes. Resolve all exceptional
fibres by successive blowups with smooth centres contained
in those fibres, writing
$\rho:\widetilde Y\to Y$ and $h=f\rho$. For each such
blowup, $\rho_*\OO_{\widetilde Y}=\OO_Y$ and
$R^j\rho_*\OO_{\widetilde Y}=0$ for $j>0$.
One local verification uses the formal neighbourhoods of
the exceptional projective bundle: their successive
conormal layers restrict to nonnegative twists
$\OO_{\PP^{r-1}}(l)$, $l\ge0$, which have no higher
cohomology. The analytic theorem on formal functions
gives the asserted vanishing; normality gives the
degree-zero equality. The Leray edge is consequently the
canonical isomorphism
\[
 R^2h_*\OO_{\widetilde Y}\xrightarrow{\simeq}
 R^2f_*\OO_Y=\mathscr G.
\tag{C.PK.5}
\]
The same identities hold for a composite by derived
pushforward.

Let $D=(h^{-1}(\Sigma))_{\mathrm{red}}$ and use the
relative logarithmic de Rham complex
\[
 \mathscr K^\bullet=
 \Omega_{\widetilde Y}^\bullet(\log D)\big/
 \bigl(h^*\Omega_{\PP^1}^1(\log\Sigma)
       \wedge\Omega_{\widetilde Y}^{\bullet-1}(\log D)\bigr).
\]
Projection onto degree zero is a morphism of complexes
$\mathscr K^\bullet\to\OO_{\widetilde Y}[0]$:
it is the identity in degree zero and zero in every
positive degree. The exact analytic comparison of D.2.3,
namely \cite[Theorem~2.18 and \S\S2.19--2.21]{Steenbrink1976},
applies to this resolution with all its fibre multiplicities
and without a K\"ahler hypothesis. It gives the logarithmic
de Rham bundle
$\mathscr H^2=\mathbb R^2h_*\mathscr K^\bullet$.
Identity monodromy makes every residue zero, by the
residue interval $[0,1)$ and the exponential comparison.
The global flat frame on $B$ therefore extends across
every point of $\Sigma$, proving
$\mathscr H^2\simeq H^2(P;\C)\otimes\OO_{\PP^1}$.
Taking degree-two direct hypercohomology of the displayed
projection and then (C.PK.5) gives a holomorphic map
\[
 p:H^2(P;\C)\otimes\OO_{\PP^1}
          \longrightarrow\mathscr G.
\tag{C.PK.6}
\]
Over $B$ it is precisely the coefficient projection
$H^2(P_c;\C)\to H^2(P_c,\OO_{P_c})$.

In an isotrivial chart this latter map is constant. It is
surjective: the class of
$\overline\alpha\wedge\overline\eta$ used in the proof
of D.1.12 is represented by a closed invariant two-form
and has nonzero Dolbeault class. The charts also show
that its kernel is a flat subspace. Choose a constant flat
class $v_0\in H^2(P;\C)$ whose projection is nonzero,
and put $s=p(v_0)\in H^0(\PP^1,\mathscr G)$.
This section is nowhere zero on $B$; zeros at points
of $\Sigma$ cause no problem.

By (C.PK.1) and the nowhere-zero section $\beta$, the
following is an exact sequence of vector bundles on the
entire projective line:
\[
 0\longrightarrow\OO_{\PP^1}\beta
  \longrightarrow\mathscr E
  \xrightarrow{\ \beta\wedge-\ }\mathscr G
  \longrightarrow0.
\tag{C.PK.7}
\]
Since $H^1(\PP^1,\OO)=0$, choose
$v\in H^0(\PP^1,\mathscr E)=H^1(Y,\OO_Y)$ with
$\beta\wedge v=s$.
We identify exactly the scalar that this lift induces in
the quotient (C.PK.2). In each isotrivial chart,
$\beta$ is the coefficient image of a fixed integral
class, $v_0$ is a fixed de Rham class, and $s=p(v_0)$
is a fixed nonzero class in $H^2(P,\OO_P)$.
The intrinsically oriented quotient $q_c$ is also fixed
by the chart: every automorphism preserves its
positive-degree class $\alpha\in\Delta^*$, so it cannot
induce inversion of $\C^*$.
Both linear maps $u\mapsto\beta\wedge u$ and
$u\mapsto\lambda(u)s$ have kernel $A$ and are nonzero.
Thus they differ by a constant $\kappa\in\C^*$ on
that chart. The constants agree on overlaps, because
all the maps and sections used in their definition are
the actual global ones. The connectedness of $B$ gives
one constant $\kappa$, and (C.PK.7) yields
\[
 \lambda_c(v_c)=\kappa^{-1}\ne0
 \quad\hbox{for every }c\in B.
\tag{C.PK.8}
\]
This proves constancy of the quotient differential,
including its normalization; it is not a boundedness
assertion about an unspecified function on $B$.

Choose $n>0$ with $nd\ge2$ and $z\in\C$ with
$\exp(2\pi\mathrm i z\kappa^{-1})=\alpha^n$.
The single ambient line $M=\exp(2\pi\mathrm i z v)$
has this quotient coordinate on every smooth fibre.
Exactly as in the $k\ne0$ case,
$M|_{P_c}=\pi_c^*L_c$ with $\deg L_c=nd\ge2$.
Its direct image has rank at least two and its relative
linear system gives a surface image over $\PP^1$,
contradicting $a(Y)=1$. This completes the proof.
\end{proof}

\begin{corollary}[C.3.11; Proved complete cohomological dichotomy]
\label{st:C.3.11}
Under the hypotheses of C.3.1, a very general smooth fibre
is a two-torus, a primary or secondary Hopf surface of
algebraic dimension zero. The Inoue exclusion is C.3.12.
Both Hopf
entries and the torus entry occur, by C.3.6 and D.2.7.
For the entire smooth locus there are exactly the following
two cohomological alternatives:
\[
\begin{array}{c|c|c}
\text{smooth fibres}&
 (h^0(\OO),h^1(\OO),h^2(\OO))&
 \text{coherent direct images}\\ \hline
\text{two-tori}&(1,2,1)&
 R^2f_*\OO_Y\simeq\det R^1f_*\OO_Y\\
b_1=1,\ b_2=0&(1,1,0)&
 R^1f_*\OO_Y\simeq\OO_{\PP^1},\
 R^2f_*\OO_Y=0 .
\end{array}
\tag{C.PK.9}
\]
The indicated cohomology dimensions hold on every
scheme-theoretic full fibre in the respective branch,
and both direct images satisfy base change there.
In the second branch
$H^q(Y,\OO_Y)=(\C,\C,0,0)$ and
$\Pic^0(Y)=\C\beta$ with the exponential convention
of D.1.9. The first line records a cohomological
alternative rather than a positivity assertion about
the splitting of $R^1f_*\OO_Y$.
\end{corollary}
\begin{proof}
C.3.7, C.3.9, C.3.10 and C.3.12 remove the bielliptic,
Kodaira and Inoue types from C.3.1; C.3.8 gives the Hopf algebraic
dimension. All smooth fibres are diffeomorphic over the
connected smooth base. In the torus branch they have
$b_1=4$, $b_2=6$, and are K\"ahler by the compact-surface
criterion. Their rational first Chern class and second
Chern number vanish by C.3.1. The surface classification
then makes them two-tori: an effective nonzero canonical
divisor has strictly positive degree against a K\"ahler
form, whereas this rational first Chern class has zero
degree. The minimal K\"ahler surfaces of Kodaira dimension
zero with these Betti numbers are exactly the two-tori.
D.1.12 supplies this branch's full-fibre assertions.

In the Hopf branch every smooth fibre has
$b_1=1$ and $b_2=0$. A nonzero holomorphic two-form
is closed and has strictly positive integral against
its conjugate, so it defines a nonzero de Rham
cohomology class. Thus $b_2=0$ forces $p_g=0$.
Since $\chi(\OO)=0$ by C.3.1, one has $q=1$.
By E4, $\mathscr E=R^1f_*\OO_Y$ is a line bundle,
and its nowhere-zero transgression section $\beta$
from D.1.9 trivializes it. The coherent sheaf
$\mathscr G=R^2f_*\OO_Y$ has rank zero, hence finite
length. Hirzebruch--Riemann--Roch on $Y$ and coherent
Leray give the exact equality
\[
 0=\chi(Y,\OO_Y)
   =\chi(\OO_{\PP^1})-\chi(\mathscr E)+\chi(\mathscr G)
   =1-1+\operatorname{length}(\mathscr G).
\]
Therefore $\mathscr G=0$. Top-degree base change and
the parameter-multiplication exact sequence give
$h^2(\OO_{Y_c})=0$ and $h^1(\OO_{Y_c})=1$ on each
full fibre; E4 gives $h^0(\OO_{Y_c})=1$.
Leray now gives the stated global structure-sheaf
cohomology. Finally $H^1(Y;\Z)=0$ and the exponential
sequence identify $\Pic^0(Y)$ with
$H^1(Y,\OO_Y)=\C\beta$, as asserted.
\end{proof}
